% ════════════════════════════════════════════════════════════════
\subsection{Kody Huffmana}
% ════════════════════════════════════════════════════════════════

Często znaki są reprezentowane przy użyciu 8 bitów (np. ASCII). Skoro częstość znaków się różni, lepiej reprezentować znaki występujące często przy użyciu krótszych ciągów bitów, a znaki występujące rzadko przy użyciu dłuższych ciągów bitów.
\begin{flalign*}
	\text{częstość} = \text{liczba wystąpień} &  &  &  &
\end{flalign*}

\paragraph{Przykład}
W pliku jest 58 znaków. Każdy znak jest elementem zbioru $\{a, e, i, s, t, r, n\}$.

\begin{table}[H]
	\centering
	\renewcommand{\arraystretch}{1.2}
	\begin{NiceTabular}{c|l|c|c|c|c|c|c|c}[margin,hvlines]
		\multicolumn{2}{c|}{}         & $a$                          & $e$ & $i$ & $s$ & $t$ & $r$ & $n$ \\
		\hline
		\multicolumn{2}{c|}{częstość} & 10                           & 15  & 12  & 3   & 4   & 13  & 1   \\
		\hline
		kod 1                         & \makecell[l]{stała długość                                       \\ słowa kodowego} & 000 & 001 & 010 & 011 & 100 & 101 & 110 \\
		\hline
		kod 2                         & \makecell[l]{zmienna długość                                     \\ słowa kodowego} & 001 & 01 & 10 & 00000 & 0001 & 11 & 00001 \\
		\hline
		kod 3                         & \makecell[l]{zmienna długość                                     \\ słowa kodowego} & 001 & 01 & 10 & 00010 & 0001 & 11 & 00001
	\end{NiceTabular}
\end{table}

Jeśli użyjemy kodu o stałej długości słowa kodowego (kod 1) to zakodowany plik ma: $58 \cdot 3 = 174$ bity. \\
Jeśli użyjemy kodu o zmiennej długości słowa kodowego (kod 2 lub kod 3) to zakodowany plik ma:
\begin{flalign*}
	10 \cdot 3 + 15 \cdot 2 + 12 \cdot 2 + 3 \cdot 5 + 4 \cdot 4 + 13 \cdot 2 + 1 \cdot 5 = 146 \text{ bitów.} &  &  &  &
\end{flalign*}

Jeśli używamy kodu o zmiennej długości słowa kodowego, to chcemy aby spełniony był warunek: \textbf{żadne słowo kodowe nie jest prefixem innego słowa kodowego}. \\
Takie kody nazywa się \uemph{kodami prefixowymi} (chociaż powinno być bezprefixowe).

Rozważmy kod 3. Ciąg $0001001$ można odkodować na dwa sposoby:
\begin{itemize}[nosep]
	\item $00010 \quad 01 \implies s \quad e$
	\item $0001 \quad 001 \implies t \quad a$
\end{itemize}
Tutaj $0001$ jest prefixem $00010$, czyli słowo kodowe znaku $t$ jest prefixem słowa kodowego znaku $s$.

\paragraph{Reprezentacja drzewowa kodów}
Słowo kodowe znaku jest ciągiem etykiet krawędzi na ścieżce od korzenia do wierzchołka zaetykietowanego tym znakiem.

\tikzset{
	mytree/.style={
			level distance=0.8cm,
			level 1/.style={sibling distance=4.0cm},
			level 2/.style={sibling distance=1.8cm},
			level 3/.style={sibling distance=1.2cm},
			level 4/.style={sibling distance=0.9cm},
			level 5/.style={sibling distance=0.7cm},
			every node/.style={circle, draw, minimum size=0.4cm, inner sep=1pt},
			elabel left/.style={draw=none, font=\scriptsize, inner sep=1pt, midway, sloped, above},
			elabel right/.style={draw=none, font=\scriptsize, inner sep=1pt, midway, sloped, above},
		}
}
% Wariant dla kroków konstrukcji: jednakowe, okrągłe węzły (etykiety typu "n,1")
\tikzset{
	htree/.style={
			mytree,
			level distance=0.95cm,
			level 4/.style={sibling distance=1.0cm},
			level 5/.style={sibling distance=0.85cm},
			every node/.style={circle, draw, minimum size=0.7cm, inner sep=0pt, font=\footnotesize, align=center},
		}
}
\begin{figure}[ht]
	\centering
	\begin{minipage}{0.48\textwidth}
		\centering
		\textbf{1. kod 1} \\ (stała długość) \\[1ex]
		\begin{tikzpicture}[mytree]
			\node {}
			child { node {}
					child { node {}
							child { node[draw=none] {a} edge from parent node[elabel left] {0} }
							child { node[draw=none] {e} edge from parent node[elabel right] {1} }
							edge from parent node[elabel left] {0}
						}
					child { node {}
							child { node[draw=none] {i} edge from parent node[elabel left] {0} }
							child { node[draw=none] {s} edge from parent node[elabel right] {1} }
							edge from parent node[elabel right] {1}
						}
					edge from parent node[elabel left] {0}
				}
			child { node {}
					child { node {}
							child { node[draw=none] {t} edge from parent node[elabel left] {0} }
							child { node[draw=none] {r} edge from parent node[elabel right] {1} }
							edge from parent node[elabel left] {0}
						}
					child { node {}
							child { node[draw=none] {n} edge from parent node[elabel left] {0} }
							child { node[draw=none] {} edge from parent[draw=none] }
							edge from parent node[elabel right] {1}
						}
					edge from parent node[elabel right] {1}
				};
		\end{tikzpicture}
	\end{minipage}
	\hfill
	\begin{minipage}{0.48\textwidth}
		\centering
		\textbf{2. kod 2} \\ (zmienna długość) \\[1ex]
		\begin{tikzpicture}[mytree]
			\node {}
			child { node {}
					child { node {}
							child { node {}
									child { node {}
											child { node[draw=none] {s} edge from parent node[elabel left] {0} }
											child { node[draw=none] {n} edge from parent node[elabel right] {1} }
											edge from parent node[elabel left] {0}
										}
									child { node[draw=none] {t} edge from parent node[elabel right] {1} }
									edge from parent node[elabel left] {0}
								}
							child { node[draw=none] {a} edge from parent node[elabel right] {1} }
							edge from parent node[elabel left] {0}
						}
					child {
							node[draw=none] {e}
							edge from parent node[elabel right] {1}
						}
					edge from parent node[elabel left] {0}
				}
			child { node {}
					child { node[draw=none] {i} edge from parent node[elabel left] {0} }
					child { node[draw=none] {r} edge from parent node[elabel right] {1} }
					edge from parent node[elabel right] {1}
				};
			% child { node[draw=none] {} edge from parent[draw=none] };
		\end{tikzpicture}
	\end{minipage}

	\vspace{1cm}
	\textbf{3. kod 3} \\ (zmienna długość) \\[1ex]
	\begin{tikzpicture}[mytree, level 4/.style={sibling distance=1.4cm}, level 5/.style={sibling distance=1.0cm}]
		\node {}
		child { node {}
				child { node {}
						child { node {}
								child { node {}
										child { node[draw=none] {} edge from parent[draw=none] }
										child { node[draw=none] {n} edge from parent node[elabel right] {1} }
										edge from parent node[elabel left] {0}
									}
								child { node {t}
										child { node[draw=none] {s} edge from parent node[elabel left] {0} }
										child { node[draw=none] {} edge from parent[draw=none] }
										edge from parent node[elabel right] {1}
									}
								edge from parent node[elabel left] {0}
							}
						child { node[draw=none] {a} edge from parent node[elabel right] {1} }
						edge from parent node[elabel left] {0}
					}
				child {
						node[draw=none] {e}
						edge from parent node[elabel right] {1}
					}
				edge from parent node[elabel left] {0}
			}
		child { node {}
				child { node[draw=none] {i} edge from parent node[elabel left] {0} }
				child { node[draw=none] {r} edge from parent node[elabel right] {1} }
				edge from parent node[elabel right] {1}
			};
		% child { node[draw=none] {} edge from parent[draw=none] };
	\end{tikzpicture}
\end{figure}

\begin{remark}
	Żadne słowo kodowe nie jest prefixem innego słowa kodowego $\iff$ wszystkie znaki są reprezentowane przez liście w reprezentacji drzewowej kodu.
\end{remark}

Etykietujemy wierzchołki w następujący sposób: każdy wierzchołek $v$ dostaje etykietę będącą sumą wystąpień znaków, które są reprezentowane przez liście poddrzewa ukorzenionego w $v$.

Niech $C$ będzie alfabetem, $f: C \to \N_+$, $f(c) =$ częstość znaku $c \in C$.
\begin{flalign*}
	d_T(c) & \text{ -- głębokość liścia reprezentującego znak } c \in C \text{ w drzewie } T                 &  & \\
	d_T(c) & = \text{długość słowa kodowego znaku } c \in C \text{ w kodzie reprezentowanym przez drzewo } T &  &
\end{flalign*}
Liczba bitów w zakodowanym pliku (kod reprezentowany przez drzewo $T$) wynosi:
\begin{flalign*}
	B(T) = \sum_{c \in C} f(c) \cdot d_T(c) &  &  &  &
\end{flalign*}

Kod prefixowy reprezentowany przez drzewo $T$ nazywamy \uemph{optymalnym}, jeśli $B(T)$ jest najmniejsze możliwe. \\
Drzewo ukorzenione nazywamy \uemph{binarnym}, jeśli każdy wierzchołek ma co najwyżej dwoje dzieci. \\
Drzewo binarne nazywamy \uemph{regularnym}, jeśli każdy wierzchołek wewnętrzny ma dokładnie dwoje dzieci.

\begin{theorem}
	Drzewo binarne reprezentujące optymalny kod jest regularne.
\end{theorem}

\begin{proof}
	Niech $T$ będzie drzewem binarnym reprezentującym optymalny kod prefixowy. Przypuśćmy, że $T$ nie jest regularne. Wtedy istnieje wierzchołek wewnętrzny $v$ z dokładnie jednym dzieckiem $w$. Niech $u$ będzie rodzicem $v$ (jeśli $v$ nie jest korzeniem $T$).

	\begin{figure}[H]
		\centering
		\begin{minipage}{0.4\textwidth}
			\centering
			\begin{tikzpicture}[mytree]
				\node {u}
				child { node {}
						child { node {}
								child { node[draw=none] {a} edge from parent node[elabel left] {0} }
								child { node[draw=none] {b} edge from parent node[elabel right] {1} }
								edge from parent node[elabel left] {0}
							}
						child { node {}
								child { node[draw=none] {c} edge from parent node[elabel left] {0} }
								child { node[draw=none] {d} edge from parent node[elabel right] {1} }
								edge from parent node[elabel right] {1}
							}
						edge from parent node[elabel left] {0}
					}
				child { node {v}
						child[missing] {}
						child { node {w}
								child { node[draw=none] {e} edge from parent node[elabel left] {0} }
								child { node[draw=none] {f} edge from parent node[elabel right] {1} }
								edge from parent node[elabel right] {1}
							}
						edge from parent node[elabel right] {1}
					};
				\node[draw=none] at (-2, 0) {$T$};
			\end{tikzpicture}
		\end{minipage}
		\hfill
		$\longrightarrow$
		\hfill
		\begin{minipage}{0.4\textwidth}
			\centering
			\begin{tikzpicture}[mytree]
				\node {u}
				child { node {}
						child { node {}
								child { node[draw=none] {a} edge from parent node[elabel left] {0} }
								child { node[draw=none] {b} edge from parent node[elabel right] {1} }
								edge from parent node[elabel left] {0}
							}
						child { node {}
								child { node[draw=none] {c} edge from parent node[elabel left] {0} }
								child { node[draw=none] {d} edge from parent node[elabel right] {1} }
								edge from parent node[elabel right] {1}
							}
						edge from parent node[elabel left] {0}
					}
				child { node {w}
						child { node[draw=none] {e} edge from parent node[elabel left] {0} }
						child { node[draw=none] {f} edge from parent node[elabel right] {1} }
						edge from parent node[elabel right] {1}
					};
				\node[draw=none] at (-2, 0) {$T'$};
			\end{tikzpicture}
		\end{minipage}
	\end{figure}

	Niech $T'$ będzie drzewem powstałym z $T$ przez usunięcie $v$ i dodanie krawędzi $uw$, jeśli $v$ nie był korzeniem $T$.
	$T'$ reprezentuje kod prefixowy i $B(T') < B(T)$ -- sprzeczność z optymalnością $T$.
\end{proof}

% ════════════════════════════════════════════════════════════════
\subsection{Algorytm Huffmana}
% ════════════════════════════════════════════════════════════════

Algorytm konstruuje drzewo reprezentujące optymalny kod prefixowy dla $(C, f)$.
Zaczynamy od $|C|$ wierzchołków izolowanych (które staną się liśćmi) i wykonujemy $|C|-1$ operacji łączenia zaczynając od najrzadszych znaków (będą miały największą głębokość w drzewie).

Będziemy używać kolejki priorytetowej z kluczem $f(c)$. Kolejkę priorytetową możemy zaimplementować używając kopca binarnego. Wtedy koszty operacji wynoszą:
\begin{itemize}[nosep]
	\item \Call{insert}{} -- $\BO(\log n)$
	\item \Call{extract\_min}{} -- $\BO(\log n)$
	\item \Call{create}{} -- $\BO(n)$
\end{itemize}
gdzie $n$ jest liczbą elementów w kolejce.

\begin{alg}{Algorytm Huffmana}{alg:huffman}
	\FunctionNN{Huffman}{$C, f$}
	\State $n \Assign |C|$
	\State $Q \Assign \Call{create}{C, f}$ \Comment{$Q$ - kolejka priorytetowa}
	\For{$i \Assign 1$ \To $n - 1$}
	\State $z \Assign \Call{new\_node}{}$
	\State $x \Assign \Call{left}{z} \Assign \Call{extract\_min}{Q}$
	\State $y \Assign \Call{right}{z} \Assign \Call{extract\_min}{Q}$
	\State $f(z) \Assign f(x) + f(y)$
	\State \Call{insert}{Q, z}
	\EndFor
	\State \Return \Call{extract\_min}{Q}
	\EndFunction
\end{alg}
Algorytm konstruuje drzewo binarne. $z$ - wierzchołek, $left, right$ - dzieci.

\begin{dygresja}
	Intuicyjnie algorytm działa zachłannie, budując drzewo \emph{od dołu}. Każdy znak $c \in C$ jest na początku osobnym, izolowanym wierzchołkiem; będzie on liściem gotowego drzewa. Kolejka priorytetowa $Q$ przechowuje korzenie wszystkich dotąd zbudowanych poddrzew, uporządkowane wg ich łącznej częstości $f$ (im rzadszy znak, tym wcześniej zostanie wyjęty).

	W każdej z $n-1$ iteracji pętli wykonujemy jeden krok łączenia:
	\begin{itemize}[nosep]
		\item \Call{extract\_min}{Q} dwukrotnie wybiera dwa korzenie o \emph{najmniejszej} częstości ($x$ oraz $y$);
		\item tworzymy nowy wierzchołek $z$ i podpinamy $x$, $y$ jako jego lewe i prawe dziecko;
		\item częstość nowego wierzchołka ustalamy jako $f(z) = f(x) + f(y)$, po czym wkładamy $z$ z powrotem do kolejki przez \Call{insert}{Q, z}.
	\end{itemize}
	Łącząc zawsze dwa najrzadsze poddrzewa, spychamy najrzadsze znaki najgłębiej w drzewie -- to one dostaną najdłuższe kody, a znaki częste pozostaną blisko korzenia i otrzymają kody krótkie. Każde łączenie zmniejsza liczbę poddrzew w $Q$ o jeden (dwa wyjęcia, jedno wstawienie), więc po $n-1$ iteracjach zostaje dokładnie jeden korzeń -- całe drzewo kodu, które zwracamy.
\end{dygresja}

\paragraph{Złożoność czasowa}
\begin{itemize}[nosep]
	\item linia 1 -- $\BO(1)$
	\item linia 2 -- $\BO(n)$
	\item wnętrze pętli z linii 3 wykona się $n - 1$ razy:
	      \begin{itemize}[nosep]
		      \item linia 4 -- $\BO(1)$
		      \item linie 5, 6 -- $\BO(\log n)$
		      \item linia 7 -- $\BO(1)$
		      \item linia 8 -- $\BO(\log n)$
	      \end{itemize}
	      Wnętrze pętli zajmuje łącznie czas $\BO(\log n)$.
\end{itemize}
\begin{flalign*}
	\BO(n) + (n - 1) \cdot \BO(\log n) = \BO(n \log n) &  &  &  &
\end{flalign*}
Złożoność czasowa algorytmu wynosi $\BO(n \log n)$.

\begin{example}[Konstrukcja drzewa Huffmana]
	Zaczynamy ze zbiorem wierzchołków: $(n, 1), (s, 3), (t, 4), (a, 10), (i, 12), (r, 13), (e, 15)$.

	% Każdy krok: para minipage obok siebie (opis | drzewo). Bloki są osobnymi
	% akapitami, więc lista łamie się między stronami. \vspace{0pt} na początku
	% minipage wymusza wyrównanie [t] do GÓRNEJ krawędzi tekstu (nie do baseline).
	\newcommand{\huffstep}[2]{%
		\noindent
		\begin{minipage}[t]{0.34\textwidth}\vspace{0pt}#1\end{minipage}%
		\hfill
		\begin{minipage}[t]{0.62\textwidth}\vspace{0pt}\centering#2\end{minipage}%
		\par\vspace{2em}%
	}

	\huffstep{%
		\textbf{1.} Łączymy $n$ i $s$. \\
		Zostają: $(t, 4), (a, 10), (i, 12), (r, 13), (e, 15)$.
	}{%
		\begin{tikzpicture}[htree]
			\node {4}
			child { node {n,1} }
			child { node {s,3} };
		\end{tikzpicture}
	}

	\huffstep{%
		\textbf{2.} Łączymy $4$ i $t$. \\
		Zostają: $(a, 10), (i, 12), (r, 13), (e, 15)$.
	}{%
		\begin{tikzpicture}[htree]
			\node {8}
			child { node {4}
					child { node {n,1} }
					child { node {s,3} }
				}
			child { node {t,4} };
		\end{tikzpicture}
	}

	\huffstep{%
		\textbf{3.} Łączymy $8$ i $a$. \\
		Zostają: $(i, 12), (r, 13), (e, 15)$.
	}{%
		\begin{tikzpicture}[htree]
			\node {18}
			child { node {8}
					child { node {4}
							child { node {n,1} }
							child { node {s,3} }
						}
					child { node {t,4} }
				}
			child { node {a,10} };
		\end{tikzpicture}
	}

	\huffstep{%
		\textbf{4.} Łączymy $i$ i $r$. \\
		Zostaje: $(e, 15)$.
	}{%
		\begin{tikzpicture}[htree]
			\node {25}
			child { node {i,12} }
			child { node {r,13} };
		\end{tikzpicture}
	}

	\huffstep{%
		\textbf{5.} Łączymy $e$ i $18$.
	}{%
		\begin{tikzpicture}[htree]
			\node {33}
			child { node {e,15} }
			child { node {18}
					child { node {8}
							child { node {4}
									child { node {n,1} }
									child { node {s,3} }
								}
							child { node {t,4} }
						}
					child { node {a,10} }
				};
		\end{tikzpicture}
	}

	\huffstep{%
		\textbf{6.} Łączymy $25$ i $33$. \\
		(ostateczne drzewo kodów)
	}{%
		\begin{tikzpicture}[htree]
			\node {58}
			child { node {25}
					child { node {i,12} edge from parent node[elabel left] {0} }
					child { node {r,13} edge from parent node[elabel right] {1} }
					edge from parent node[elabel left] {0}
				}
			child { node {33}
					child { node {e,15} edge from parent node[elabel left] {0} }
					child { node {18}
							child { node {8}
									child { node {4}
											child { node {n,1} edge from parent node[elabel left] {0} }
											child { node {s,3} edge from parent node[elabel right] {1} }
											edge from parent node[elabel left] {0}
										}
									child { node {t,4} edge from parent node[elabel right] {1} }
									edge from parent node[elabel left] {0}
								}
							child { node {a,10} edge from parent node[elabel right] {1} }
							edge from parent node[elabel right] {1}
						}
					edge from parent node[elabel right] {1}
				};
		\end{tikzpicture}
	}

	Odczytany kod z drzewa w kroku 6:
	\begin{flalign*}
		i & \implies 00    &  & \\
		r & \implies 01    &  & \\
		e & \implies 10    &  & \\
		n & \implies 11000 &  & \\
		s & \implies 11001 &  & \\
		t & \implies 1101  &  & \\
		a & \implies 111   &  &
	\end{flalign*}
\end{example}
